\addtocontents{toc}{\protect\Needspace{3\baselineskip}}
\originalchapter{大 N 展开与弦的拓扑}\label{ch:largeN}
\originalsection{双线图计数}
对 adjoint 矩阵场, 取 't Hooft 耦合 $\lambda=g_{\rm YM}^2N$ 固定. 把作用量写为 $N\operatorname{Tr}[\cdots]/\lambda$, 每传播子带 $1/N$, 每顶点带 $N$, 每个闭合颜色指标回路带 $N$. 于是连通真空图的 $N$ 次方为 $N^{F-E+V}$. 双线图厚化为可定向曲面, Euler 特征数 $F-E+V=2-2h$, 所以亏格 $h$ 对应 $N^{2-2h}$.

\begin{center}\begin{tikzpicture}[scale=.8]
\draw[very thick,structurecolor](-2,-.15)--(2,-.15);
\draw[very thick,structurecolor](-2,.15)--(2,.15);
\draw[-{Stealth}](-.8,.15)--(.8,.15);
\draw[-{Stealth}](.8,-.15)--(-.8,-.15);
\node at(0,.75){矩阵传播子: 两条颜色指标线};
\node at(-2.4,.25){$i$};\node at(-2.4,-.25){$j$};
\node at(2.4,.25){$k$};\node at(2.4,-.25){$l$};
\end{tikzpicture}\end{center}
图中双线是指标流向, 不是两条不同的物理弦. 对带边界的图, fundamental 回路引入边界数 $b$, 权重 $N^{2-2h-b}$. 这与弦世界面拓扑权重相似, 可提示 $g_s\sim1/N$, 但没有自动构造目标空间几何或局域世界面作用量.

\originalsection{因子化与局限}
规范不变单迹算符按适当 $N$ 归一化后, 连接两算符的图比乘积的断开图少两个 $N$ 次方, 故 $\langle\mathcal O_1\mathcal O_2\rangle-\langle\mathcal O_1\rangle\langle\mathcal O_2\rangle=O(N^{-2})$. 这类似经典变量的涨落被抑制, 但颜色平面图仍含所有耦合阶次; $N\to\infty$ 不等于 $\lambda\to0$.

Wilson 环为闭路径的平行输运迹. 未归一化 fundamental 与 adjoint 迹满足逐个群元素的恒等式 $\operatorname{Tr}_{\rm adj}U=|\operatorname{Tr}_{F}U|^2-1$. 大 $N$ 因子化使 adjoint 面积律的领先指数约为 fundamental 的两倍. 有限 $N$ 时 adjoint 电荷可以被胶子屏蔽, 所以极长距离的结论依取极限顺序而变.

\originalsection{一次明确的拓扑计数}
考虑矩阵标量理论的四次顶点真空图. 单顶点把两对腿收缩成两条传播子, 有 $V=1,E=2$. 在平面收缩中三个独立指标面 $F=3$, 权重 $N^{1-2+3}=N^2$, Euler 特征数为 2, 对应球面. 对交叉收缩则只有一个指标面 $F=1$, 权重 $N^{1-2+1}=N^0$, Euler 特征数为 0, 对应一个柄. 两图的动量积分可以不同, 但颜色压低 $N^{-2}$ 来自指标拓扑, 不依靠某一次动量积分巧合.

这项计数也说明 leading planar 并不等于仅保留很少几个 Feynman 图. 每个 $\lambda$ 阶次可以有许多平面图, 强耦合时它们须整体求和. 弦几何若能重组这些图, 所解决的正是这一困难, 而不是把 $N$ 很大误当作所有相互作用都弱.
